% Tabela: Delta de Cliff para LCP — comparações par a par selecionadas
\begin{table*}[htbp]
\centering
\footnotesize
\caption{\novo{Delta de Cliff ($\delta$) para LCP — comparações par a par selecionadas por tamanho de conjunto de dados}}
\begin{tabular}{|l|c|c|c|c|c|c|}
\hline
\textbf{Comparação} & \textbf{XS} & \textbf{S} & \textbf{M} & \textbf{L} & \textbf{XL} & \textbf{Interpretação} \\
\hline
CSR (base) vs.\ CSR + Gzip           & 1,00 & 1,00 & 1,00 & 1,00 & 1,00 & Grande \\
CSR (base) vs.\ CSR + Minificação    & 1,00 & 1,00 & 1,00 & 1,00 & 1,00 & Grande \\
CSR + Gzip vs.\ CSR + Min.\ + Gzip  & 0,66 & 1,00 & 1,00 & 1,00 & 1,00 & Grande \\
CSR + Min.\ + Gzip vs.\ SSR + Gzip  & 1,00 & 1,00 & 1,00 & 1,00 & 1,00 & Grande \\
SSR (base) vs.\ SSR + Gzip          & 1,00 & 1,00 & 1,00 & 1,00 & 1,00 & Grande \\
\hline
\end{tabular}
\vspace{2pt}
\begin{flushleft}
\footnotesize
\textbf{*} XS = Extra-Small; S = Small; M = Medium; L = Large; XL = Extra-Large.
$\delta = 1{,}0$ com IC$_{95\%}$ $[1{,}0;\,1{,}0]$ indica separação completa entre grupos.
Escala: $|\delta| < 0{,}147$ = negligível; $0{,}147$--$0{,}330$ = pequeno; $0{,}330$--$0{,}474$ = médio; $|\delta| \geq 0{,}474$ = grande~\cite{Romano2006}.
IC de 95\% calculados via bootstrap~\cite{efron:1994}.
\end{flushleft}
\label{tab:cliffs-delta}
\end{table*}
